% Part: many-valued-logic
% Chapter: three-valued-logics
% Section: introduction

\documentclass[../../../include/open-logic-section]{subfiles}

\begin{document}

\olfileid{mvl}{thr}{int}

\olsection{ভূমিকা}

$\True$ ও~$\False$-এর সঙ্গে আর একটি মান~$\Undef$ যোগ করলেই আমরা একটি
ত্রিমানী যুক্তিবিদ্যা পাই। সত্যমান মাত্র একটি বাড়লেও $\lnot$,
$\land$, $\lor$ ও~$\lif$-এর সত্যমান-অপেক্ষক সংজ্ঞায়িত করার সম্ভাব্য
উপায়ের সংখ্যা অনেক। তখন কোনো যুক্তিবিদ্যা এই সত্যমান-অপেক্ষকগুলির
যেকোনো সমন্বয় ব্যবহার করতে পারে; আর শুধু~$\True$-কে, অথবা $\True$
ও~$\Undef$ উভয়কেই, মনোনীত করার বিকল্পও থাকে।

এখানে আমরা সর্বাধিক পরিচিত কয়েকটি ত্রিমানী যুক্তিবিদ্যা, সেগুলির
প্রেরণা এবং কয়েকটি বৈশিষ্ট্য উপস্থাপন করব।

\end{document}
